\section{Discussion}
\label{sec:discussion}

\paragraph{An LLM judge is necessary, not optional.} The most actionable methodological lesson is
negative: the standard discriminative tool for scientific claim verification---NLI entailment over
claim--evidence pairs \citep{wadden2020fact}---is not a safe stand-alone measure of whether science
supports a policy claim. It systematically labels genuine support as ``neutral,'' because an
abstract seldom restates a press-release claim in entailing language, and the resulting downward
bias is large enough to reverse a substantive conclusion (here, the apparent partisan gap in
evidence quality). A high-capacity generative judge, which reasons about whether the study's
findings \emph{bear on} the claim rather than whether the texts \emph{entail} one another, is
required to recover the right answer. This does not make NLI useless: it is cheap, fast, and
runs on every pair, which is exactly what a surrogate needs to be.

\paragraph{Validity without humans.} The combination that makes the pipeline both autonomous and
valid is surrogate-plus-judge-plus-PPI. Prediction-powered inference \citep{angelopoulos2023ppi}
and design-based supervised learning \citep{egami2023dsl} were developed to let researchers use
abundant imperfect labels alongside a small set of trusted labels; our contribution is to observe
that the ``trusted'' labels need not come from humans. A frontier LLM judge supplies
gold-equivalent labels at the scale of a few dozen to a few hundred items---enough to estimate and
correct the surrogate's bias---while the open-source NLI models supply the breadth. The
human research assistant, the traditional bottleneck and cost center of this kind of measurement,
is removed without sacrificing the inferential guarantees that justify the estimate.

\paragraph{Diagnose across methods, not within.} A practical corollary concerns how to know
whether an automated coder can be trusted. High agreement \emph{among} models of the same family
is reassuring but can be an illusion of correctness: our two NLI models agreed at $\kappa=0.65$
while sharing the very bias that distorted the estimate. The informative comparison is
\emph{across} paradigms---generative versus discriminative---and the alternative-annotator test
\citep{calderon2025alttest} formalizes when one may substitute for the other. We recommend reporting
cross-method agreement as a routine diagnostic in LLM-based measurement, precisely because
within-method agreement hid the problem here.

\paragraph{What we learn about AI policymaking.} Substantively, two findings stand out. First,
evidence use is higher than a naive automated reading suggests: most technical claims California
legislators make about AI are supported by the peer-reviewed literature our pipeline retrieves
(roughly 0.6--0.7 for both parties). Whatever one's priors about ``post-truth'' politics, in this
domain and at the level of discrete factual claims, legislators' assertions are largely backed by
science that exists and is mostly open-access. Second, the robust partisan asymmetry is in
\emph{engagement}, not \emph{accuracy}: Democrats talk about AI far more, and more intensively, than
Republicans, but once a competent judge is used the two parties' claims are supported at similar
rates. The headline ``Democrats use better evidence'' implied by the naive measure is an artifact of
the instrument, a cautionary tale for the rapidly growing practice of using off-the-shelf NLP to
score political text.

\paragraph{A template for computational social science.} Beyond this application, the design is a
reusable template: pair a cheap, scalable open-source classifier (the surrogate) with a small,
high-quality LLM-judge gold set, reconcile them with PPI/DSL, and validate by cross-method
agreement. It runs with no proprietary API and no human coders, which lowers both the cost and the
access barriers to large-scale, valid text measurement---and makes the whole analysis exactly
reproducible from released code and data.
